\documentclass{article}

% if you need to pass options to natbib, use, e.g.:
%     \PassOptionsToPackage{numbers, compress}{natbib}
% before loading neurips_2026

% The authors should use one of these tracks.
% Before accepting by the NeurIPS conference, select one of the options below.
% 0. "default" for submission
\usepackage{neurips_2026}
% the "default" option is equal to the "main" option, which is used for the Main Track with double-blind reviewing.
% 1. "main" option is used for the Main Track
%  \usepackage[main]{neurips_2026}
% 2. "position" option is used for the Position Paper Track
%  \usepackage[position]{neurips_2026}
% 3. "eandd" option is used for the Evaluations & Datasets Track
 % \usepackage[eandd]{neurips_2026}
 % if you need to opt-in for a single-blind submission in the E&D track:
 %\usepackage[eandd, nonanonymous]{neurips_2026}
% 4. "creativeai" option is used for the Creative AI Track
%  \usepackage[creativeai]{neurips_2026}
% 5. "sglblindworkshop" option is used for the Workshop with single-blind reviewing
 % \usepackage[sglblindworkshop]{neurips_2026}
% 6. "dblblindworkshop" option is used for the Workshop with double-blind reviewing
%  \usepackage[dblblindworkshop]{neurips_2026}

% After being accepted, the authors should add "final" behind the track to compile a camera-ready version.
% 1. Main Track
 % \usepackage[main, final]{neurips_2026}
% 2. Position Paper Track
%  \usepackage[position, final]{neurips_2026}
% 3. Evaluations & Datasets Track
 % \usepackage[eandd, final]{neurips_2026}
% 4. Creative AI Track
%  \usepackage[creativeai, final]{neurips_2026}
% 5. Workshop with single-blind reviewing
%  \usepackage[sglblindworkshop, final]{neurips_2026}
% 6. Workshop with double-blind reviewing
%  \usepackage[dblblindworkshop, final]{neurips_2026}
% Note. For the workshop paper template, both \title{} and \workshoptitle{} are required, with the former indicating the paper title shown in the title and the latter indicating the workshop title displayed in the footnote.
% For workshops (5., 6.), the authors should add the name of the workshop, "\workshoptitle" command is used to set the workshop title.
% \workshoptitle{WORKSHOP TITLE}

% "preprint" option is used for arXiv or other preprint submissions
 % \usepackage[preprint]{neurips_2026}

% to avoid loading the natbib package, add option nonatbib:
%    \usepackage[nonatbib]{neurips_2026}

\usepackage[utf8]{inputenc} % allow utf-8 input
\usepackage[T1]{fontenc}    % use 8-bit T1 fonts
\usepackage{hyperref}       % hyperlinks
\usepackage{url}            % simple URL typesetting
\usepackage{booktabs}       % professional-quality tables
\usepackage{amsfonts}       % blackboard math symbols
\usepackage{amsmath,amssymb}
\usepackage{nicefrac}       % compact symbols for 1/2, etc.
\usepackage{microtype}      % microtypography
\usepackage{xcolor}         % colors
\usepackage{graphicx}
\usepackage{multirow}
\usepackage{array}
\usepackage{siunitx}
\usepackage{caption}
\usepackage{subcaption}
\usepackage{listings}
\usepackage{algorithm}
\usepackage{algpseudocode}

\lstdefinestyle{cuda}{
  basicstyle=\ttfamily\footnotesize,
  keywordstyle=\color{blue!70!black}\bfseries,
  commentstyle=\color{green!50!black}\itshape,
  stringstyle=\color{red!70!black},
  showstringspaces=false,
  breaklines=true,
  numbers=left,
  numberstyle=\tiny\color{gray},
  frame=single,
  language=C++,
}

\title{AccelEval: Evaluating and Decomposing LLM-Generated CUDA Code at Scale}


% The \author macro works with any number of authors. There are two commands
% used to separate the names and addresses of multiple authors: \And and \AND.
%
% Using \And between authors leaves it to LaTeX to determine where to break the
% lines. Using \AND forces a line break at that point. So, if LaTeX puts 3 of 4
% authors names on the first line, and the last on the second line, try using
% \AND instead of \And before the third author name.


\author{%
  Chen Dong \\
  % Affiliation \\
  \texttt{zijieteam@gmail.com} \\
  % \And
  % Co-Author \\
  % Affiliation \\
  % \texttt{email} \\
}


\begin{document}


\maketitle


\begin{abstract}
We introduce \textsc{AccelEval}, a benchmark for evaluating large language models (LLMs) on translating sequential CPU programs into optimized CUDA kernels. Unlike prior kernel-generation benchmarks that focus on isolated GPU primitives or PyTorch module replacement, \textsc{AccelEval} stresses \emph{end-to-end} acceleration across \textbf{43 production-style workloads} spanning scientific computing, graph analytics, financial engineering, dynamic programming, and operations research. Each task is compiled and executed as a whole program against a deterministic validation oracle, and is wall-clock timed from host pointer to host result---closing a loophole in prior setups where computation can be hidden in untimed initialization. We evaluate \textbf{seven state-of-the-art models} (Gemini~3.1~Pro, Claude~Opus~4.6, GPT-5.4, Qwen~3.6~Plus, Kimi~K2.5, GLM~5.1, DeepSeek~V3.2) at three input scales under a zero-hint prompt (L3). On the medium scale---our leaderboard setting---the strongest model (Gemini~3.1~Pro) passes 34 of 42 tasks, attains a geometric-mean speedup of $30.9\times$, and reaches $\mathrm{fast}_{10}=61.9\%$, while other models cluster between 24 and 33 passing tasks (57\% to 79\%). Beyond raw scoring, we formulate CUDA optimization as a \emph{strategy recommendation} problem over a structured algorithm space and release three calibrated baselines (nearest-neighbor transfer, greedy coordinate search, linear regression) plus a pairwise-diff attribution pipeline that identifies which concrete optimization patterns drive each speedup jump.
\end{abstract}



\section{Introduction}
\label{sec:intro}

GPU acceleration is central to modern scientific computing, machine learning, financial engineering, graph analytics, and large-scale optimization. Yet writing efficient CUDA remains difficult: programmers must reason about parallelism, memory hierarchy, host--device transfer, synchronization, numerical stability, and hardware-specific libraries. As large language models (LLMs) become increasingly capable code generators, a natural question is: \emph{can they automate CPU-to-CUDA acceleration end-to-end, producing code that is not only correct but also practically faster?}

Existing benchmarks only partially answer this question. General code benchmarks such as HumanEval~\citep{humaneval}, MBPP~\citep{mbpp}, SWE-bench~\citep{swebench}, LiveCodeBench~\citep{livecodebench}, and BigCodeBench~\citep{bigcodebench} measure functional correctness, but do not evaluate performance-critical GPU programming. GPU-focused benchmarks such as KernelBench~\citep{kernelbench}, TritonBench~\citep{tritonbench}, ComputeEval~\citep{computeeval}, ParEval~\citep{pareval}, and MultiKernelBench~\citep{multikernelbench} move closer to this setting, but typically focus on isolated kernels, short snippets, or framework-level operator replacement. They therefore leave open the harder question of whole-program acceleration: starting from a sequential CPU reference and producing a complete CUDA implementation that handles memory management, host--device transfer, kernel orchestration, and validated output at realistic input scales.

\textsc{AccelEval} is designed around this end-to-end view. Each task exposes a single function, \texttt{solution\_compute}, that receives host-side inputs and must perform all device allocation, H2D transfer, kernel or library execution, D2H transfer, and cleanup inside the timed region. This design closes timing loopholes that arise when initialization or data movement is outside the measured path; in early versions of our harness, one frontier model achieved an apparent $10^{6}\times$ speedup by moving computation into an untimed setup routine. Each solution is checked against a deterministic CPU oracle and evaluated at three input scales, allowing us to separate correctness from scaling behavior.

Beyond leaderboard results, we study \emph{why} generated CUDA code is fast. We formulate optimization as a strategy-recommendation problem over a structured algorithm space, introduce classical baselines for this setting, and build a pairwise-diff attribution pipeline that connects speedup improvements to specific optimization patterns such as tiling, warp-level reduction, template dispatch, library calls, and problem-specific precomputation.

\paragraph{Contributions.}
\begin{enumerate}
  \item We release \textsc{AccelEval}, a curated suite of \textbf{43 CPU-to-CUDA tasks} across nine categories, each with a CPU reference, a deterministic validation oracle, and three input scales ($\sim$9\,ms / $\sim$0.25\,s / $\sim$5\,s median CPU baseline).
  \item We adopt a \emph{unified single-function interface} (\texttt{solution\_compute}) that is evaluated end-to-end, closing the initialization and harness-level loopholes that can arise when only an inner kernel is timed.
  \item We present a large-scale L3 evaluation of \textbf{seven} frontier LLMs at three input scales. On the medium leaderboard, Gemini~3.1~Pro leads with 34/42 passing tasks and a $30.9\times$ geometric-mean speedup; the runners-up trail by 1--10 passing tasks and $2$--$8\times$ in geometric-mean speedup.
  \item We formulate CUDA optimization as strategy recommendation in a structured algorithm space $\mathcal{A}_0$, and release three calibrated baselines (nearest-neighbor transfer, greedy coordinate search, linear regression) to anchor what non-LLM methods can achieve from the same experience pool.
  \item We introduce a \emph{decomposition pipeline}: an auto-detected optimization pattern database combined with pairwise-diff attribution over multi-turn agent runs, enabling causal claims of the form ``pattern $P$ contributes $\Delta$ speedup on problem class $C$.''
\end{enumerate}



\section{Related Work}
\label{sec:related}

\paragraph{General code-generation benchmarks.}
HumanEval~\citep{humaneval} and MBPP~\citep{mbpp} established unit-test-based evaluation for short Python functions. Later benchmarks broadened the setting: SWE-bench~\citep{swebench} evaluates repository-level issue resolution, LiveCodeBench~\citep{livecodebench} emphasizes contamination-resistant competitive programming, and BigCodeBench~\citep{bigcodebench} tests complex library use and instruction following. These benchmarks are useful for measuring general coding ability, but they do not target GPU programming or runtime performance.

\paragraph{GPU kernel benchmarks.}
KernelBench~\citep{kernelbench} is the closest prior benchmark: it asks models to replace PyTorch module bodies with custom CUDA or Triton kernels and scores speedup relative to framework baselines. TritonBench~\citep{tritonbench} focuses on Triton operator generation, while MultiKernelBench~\citep{multikernelbench} extends kernel generation to multiple accelerator backends. ComputeEval~\citep{computeeval} tests CUDA code generation on short prompts, and ParEval~\citep{pareval} studies parallel code generation across CUDA, OpenMP, MPI, HIP, and Kokkos. These benchmarks show that performance-aware code generation is distinct from ordinary functional synthesis, but they share three limitations that \textsc{AccelEval} is designed to address:

\begin{enumerate}
  \item \emph{Granularity.} Prior benchmarks typically score a single kernel or short snippet rather than a complete accelerated program. Host-side memory management, multi-kernel orchestration, persistent device buffers, and library dispatch are often hidden by the harness or outside the task scope.
  \item \emph{Timing boundaries.} Timing only an inner kernel can exclude initialization, data movement, or precomputation. In early versions of our harness, one frontier model obtained an apparent $10^{6}\times$ speedup by moving computation into an untimed setup routine. \textsc{AccelEval}'s single \texttt{solution\_compute} entry point places the full host-to-host path under the benchmark contract.
  \item \emph{Input scale.} Most prior suites evaluate a single input size. This can obscure scaling behavior: a GPU implementation that is slower than CPU on small inputs may become much faster at scale, while a naive parallelization may pass small tests but fail to exploit larger workloads. \textsc{AccelEval} reports three deterministic input scales per task.
\end{enumerate}

Table~\ref{tab:bench-compare} summarizes these axes across recent suites.

\begin{table}[ht]
\centering
\small
\caption{Comparison of LLM code benchmarks on axes relevant to GPU acceleration. ``Scope'' is the unit of generation; ``timed'' indicates whether runtime speedup is measured; ``end-to-end'' means host-to-device transfer, computation, and device-to-host transfer are included in the benchmark contract; ``scales'' is the number of input sizes reported per task.}
\label{tab:bench-compare}
\begin{tabular}{@{}lllccc@{}}
\toprule
\textbf{Benchmark} & \textbf{Scope} & \textbf{Target} & \textbf{Timed} & \textbf{End-to-end} & \textbf{Scales} \\
\midrule
HumanEval      & function   & Python          & --- & ---  & 1 \\
MBPP           & function   & Python          & --- & ---  & 1 \\
SWE-bench      & repo patch & Python          & --- & ---  & 1 \\
ComputeEval    & snippet    & CUDA C++        & --- & ---  & 1 \\
KernelBench    & single kernel & CUDA/Triton  & \checkmark & partial & 1 \\
TritonBench    & single kernel & Triton       & \checkmark & partial & 1 \\
ParEval        & snippet    & CUDA/OpenMP/MPI & \checkmark & partial & 1 \\
MultiKernelBench & single kernel & CUDA/ROCm/Ascend & \checkmark & partial & 1 \\
BabelTower     & translation & C$\to$CUDA     & --- & --- & 1 \\
\midrule
\textsc{AccelEval} (ours) & whole program & CPU C $\to$ CUDA & \checkmark & \checkmark & 3 \\
\bottomrule
\end{tabular}
\end{table}

\paragraph{Auto-parallelization and program translation.}
Classical source-to-source systems such as OpenACC, Bones~\citep{bones}, Par4All~\citep{par4all}, and PPCG~\citep{ppcg} translate analyzable loop nests to parallel code using directives, skeletons, or polyhedral analysis. More recent learning-based systems, including BabelTower~\citep{babeltower} and CodeRosetta~\citep{coderosetta}, learn program translation or parallelization patterns from code corpora. These systems are closely related to CPU-to-GPU translation, but they are usually evaluated by translation quality, compilation, or functional correctness rather than end-to-end speedup against a fixed CPU reference.

\paragraph{HPC benchmarks and auto-tuning.}
The tasks in \textsc{AccelEval} draw on established HPC benchmark suites and mini-apps, including GAP~\citep{gapbs}, HPCG~\citep{hpcg}, NAS Parallel Benchmarks~\citep{npb}, Rodinia~\citep{rodinia}, and miniWeather~\citep{miniweather}. We repurpose their computational kernels as LLM acceleration targets under a uniform correctness and timing harness. Our strategy-recommendation analysis (Section~\ref{sec:formulation}) is also related to auto-tuning systems such as ATLAS~\citep{atlas}, OpenTuner~\citep{opentuner}, Kernel Tuner~\citep{kerneltuner}, and learned schedule search in Halide and TVM~\citep{halide-autoscheduler,tvm}. Unlike these systems, we study how a natural-language model selects and implements optimization strategies, and we use classical search and regression methods (Section~\ref{sec:strategy}) as non-LLM baselines.

\section{Problem Formulation}
\label{sec:formulation}

We cast CPU-to-CUDA optimization as a \emph{strategy recommendation}
problem over a structured algorithm space. This framing lets us study
LLMs not only as black-box generators, but as policies operating in a
well-defined decision space, enabling both principled baselines and
decomposition of where their advantage comes from.

\subsection{Problem Space $\mathcal{P}$}
Let $\mathcal{P} = \{Q_1, \ldots, Q_n\}$ be a set of CPU reference
snippets (one per \textsc{AccelEval} task). An \emph{affinity function}
$\rho : \mathcal{P} \times \mathcal{P} \to \mathbb{R}_{\geq 0}$
measures structural similarity, making $(\mathcal{P}, \rho)$ a weighted
graph. In our instantiation, $\rho$ is induced by a kernel over
static-analysis features (Section~\ref{sec:feature}).

\subsection{Feature Space $\mathcal{F}$}
\label{sec:feature}
A feature extractor $f : \mathcal{P} \to \mathcal{F} \subseteq
\mathbb{R}^d$ maps each problem to a vector
$\mathbf{q} = f(Q)$ capturing loop-nest depth, stride-1 access ratio,
arithmetic intensity, indirection level, reduction presence, etc.
Affinity is realized via a kernel:
\begin{equation}
\rho(Q_i, Q_j) = \kappa_\mathcal{F}\bigl(f(Q_i), f(Q_j)\bigr).
\label{eq:affinity}
\end{equation}

\subsection{Algorithm Space $\mathcal{A}$}
Let $\mathcal{T} = \{t_1, \ldots, t_m\}$ be a catalog of
transformation rules (Table~\ref{tab:trans}). Each rule $t_i$ carries
an intensity level $s_i \in \{0, 1, \ldots, M_i\}$, with $s_i = 0$
meaning inactive. A \emph{strategy vector} is
\begin{equation}
\mathbf{s} = (s_1, \ldots, s_m) \in \mathcal{A}
   = \prod_{i=1}^{m}\{0, 1, \ldots, M_i\}.
\label{eq:strategy}
\end{equation}

\begin{table}[ht]
\centering
\small
\caption{Example transformation rules with intensity semantics.}
\label{tab:trans}
\begin{tabular}{@{}clll@{}}
\toprule
$s_i$ & \textbf{Rule} & \textbf{Intensity semantics} & $M_i$ \\
\midrule
$s_1$ & Shared-memory tiling & $0$: off; $\ell$: tile level $\ell$ & 4 \\
$s_2$ & Coalescing layout     & $0$: none; $1$: row-major; $2$: col-major; $3$: space-filling & 3 \\
$s_3$ & Loop unroll           & $0$: off; $\ell$: factor $2^{\ell}$ & 5 \\
$s_4$ & Reduction strategy    & $0$: scalar; $1$: warp shuffle; $2$: shared-tree; $3$: atomic & 3 \\
$s_5$ & Loop fusion           & $0$: none; $\ell$: fuse $\ell{+}1$ adjacent loops & 3 \\
\bottomrule
\end{tabular}
\end{table}

Rules interact through prerequisites ($\mathcal{E}_{\mathrm{dep}}$),
mutual exclusion ($\mathcal{E}_{\mathrm{excl}}$), and
problem-dependent intensity caps $\bar{M}_i(\mathbf{q})$, yielding the
feasible set
\begin{multline}
\mathcal{A}_0 = \Bigl\{\,
  \mathbf{s} \in \mathcal{A} \mathrel{\Big|}
  \bigwedge_{(i,j)\in \mathcal{E}_{\mathrm{dep}}}
    \bigl(s_j \geq 1 \Rightarrow s_i \geq 1\bigr) \;\land \\
  \bigwedge_{(i,j)\in \mathcal{E}_{\mathrm{excl}}}
    \min(s_i, s_j) = 0 \;\land\;
  \bigwedge_{i=1}^{m} s_i \leq \bar{M}_i(\mathbf{q})
\,\Bigr\}.
\label{eq:feasible}
\end{multline}
A strategy kernel distinguishes ordered from categorical dimensions:
\begin{equation}
\kappa_\mathcal{A}(\mathbf{s},\mathbf{s}')
  = \exp\!\Biggl(-\sum_{i=1}^{m} w_i \, d_i(s_i, s_i')\Biggr),\;\;
d_i(s_i, s_i') = \begin{cases}
|s_i - s_i'|/M_i, & \text{ordered},\\
\mathbb{1}[s_i \neq s_i'], & \text{categorical}.
\end{cases}
\label{eq:skernel}
\end{equation}

\subsection{Performance Space $\mathcal{Y}$}
The reward is the achieved end-to-end speedup, clipped to zero on
failure:
\begin{equation}
y(\mathbf{s}, Q) = \frac{T_{\mathrm{CPU}}(Q)}{T_{\mathrm{GPU}}\bigl(\mathrm{Apply}(Q, \mathbf{s})\bigr)},
\qquad
y = 0 \text{ if compilation fails or output is incorrect}.
\label{eq:reward}
\end{equation}

\subsection{Observation Model}
Each multi-turn experiment yields a tuple
$E_l = (Q_{j(l)}, \mathbf{s}_l^A, \mathbf{s}_l^B, r_l)$, where
$\mathbf{s}^A$ and $\mathbf{s}^B$ are the strategy before and after
an edit. Define the strategy delta
$\Delta \mathbf{s}_l = \mathbf{s}_l^B - \mathbf{s}_l^A$.
We assume a noisy \emph{incremental effect} model:
\begin{equation}
r_l = g\bigl(f(Q_{j(l)}),\; \mathbf{s}_l^A,\; \Delta \mathbf{s}_l\bigr)
      + \epsilon_l, \qquad
\epsilon_l \sim \mathcal{N}(0, \sigma^2).
\label{eq:incremental}
\end{equation}
The three-argument form of $g$ captures that the marginal effect of a
change $\Delta\mathbf{s}$ depends on the current baseline
$\mathbf{s}^A$: \emph{e.g.}, further unroll on top of already-tiled code
behaves differently than unroll on untiled code.

\subsection{Problem Statement}
\label{sec:problem}
\begin{quote}
\textbf{Problem (Strategy Recommendation).}
Given the problem space $\mathcal{P}$ with feature extractor $f$,
the algorithm space $\mathcal{A}$ with feasible subset
$\mathcal{A}_0$, and a history of experiments
$\{E_l\}_{l=1}^{k}$, a new query
$(Q^*, \mathbf{s}^A)$ with feature
$\mathbf{q}^* = f(Q^*)$ (and $\mathbf{s}^A = \mathbf{0}$ when starting
from raw CPU), solve
\begin{equation}
\mathbf{s}^* = \arg\max_{\mathbf{s} \in \mathcal{A}_0}
  \hat{g}\bigl(\mathbf{q}^*,\; \mathbf{s}^A,\; \mathbf{s} - \mathbf{s}^A\bigr),
\label{eq:objective}
\end{equation}
where $\hat{g}$ is the effect-function estimator learned from history.
\end{quote}

This formulation unifies three different views of the problem: (i)~a
\emph{retrieval} view (pick a strategy that worked on a similar
problem); (ii)~a \emph{combinatorial-search} view (greedily optimize
each rule coordinate); and (iii)~a \emph{regression} view (fit
$\hat{g}$ and maximize). Section~\ref{sec:strategy} contrasts the
three.

\section{Benchmark Design}
\label{sec:design}

\textsc{AccelEval} consists of 43 tasks curated from established
high-performance computing suites, domain-specific scientific codes,
classical algorithm texts, and industrial operations-research workloads.
For each task, we extract the algorithmically meaningful kernel, rewrite
it against a \emph{single unified interface}, and ship three deterministic
input scales (small, medium, large). This section describes each of these
steps.

\subsection{Task Sources and Extraction}
\label{sec:sources}

\paragraph{Source suites.}
The 43 tasks are drawn from five families of prior work.

\begin{itemize}
  \item \textbf{HPC reference benchmarks} — kernels from the GAP Benchmark
  Suite~\citep{gapbs} (connected components, PageRank, triangle
  counting), HPCG~\citep{hpcg} (sparse matrix--vector, symmetric
  Gauss--Seidel, multigrid V-cycle), the NAS Parallel Benchmarks~\citep{npb}
  (conjugate-gradient, LU-SSOR, ADI pentadiagonal), the Rodinia
  suite~\citep{rodinia} (BFS levels), and the
  miniWeather~\citep{miniweather} mini-app.
  \item \textbf{Scientific / simulation codes} — Smoothed Particle
  Hydrodynamics (SPH) kernels for cell indexing, force computation, and
  position integration; a heat-diffusion stencil; a discretized Hamilton--Jacobi--Bellman
  solver for dynamic pricing.
  \item \textbf{Classical algorithms} — textbook problems with well-known
  complexity: single-source shortest paths (Bellman--Ford), maximum flow
  (push--relabel), held-Karp traveling salesman, Smith--Waterman sequence
  alignment, dynamic-time warping, DBSCAN clustering, Euclidean and
  Hausdorff distance matrices, and Thompson NFA regex matching.
  \item \textbf{Financial computing} — options pricing kernels (Black--Scholes,
  Monte-Carlo path simulation, bond and repo pricing) and a batched
  low-rank principal component analysis for portfolio construction, based
  on FinanceBench-style workloads.
  \item \textbf{Operations research} — dynamic-programming formulations for
  network revenue management, inventory replenishment, Gittins-index
  computation, robust value iteration, self-exciting pricing, Nash flows
  over time, along with a first-order linear-programming solver and a
  crew-pairing instance.
\end{itemize}

\paragraph{Extraction protocol.}
For each source program we apply a three-step extraction. First, we
\emph{isolate} the numerically dominant region by profiling the original
code and selecting the region that consumes over 90\% of the runtime.
Second, we \emph{decouple} that region from any surrounding I/O, driver
harness, or library dependency; what remains is a pure-computation kernel
expressed as a handful of C functions operating over host-visible buffers.
Third, we \emph{rewrap} the kernel behind the unified interface of
Section~\ref{sec:interface}, so that each benchmark, regardless of
pedigree, exposes exactly the same contract to the language model.

This pipeline removes two confounds: (i) legacy suites often couple the
kernel to bespoke input parsers or build systems, making cross-task
comparison meaningless; (ii) external library dependencies (BLAS,
pthreads, MPI) obscure what the GPU code must replicate. After extraction,
every task has the same shape: deterministic binary inputs, a single C
reference, and a well-specified output.

\subsection{End-to-End Timing and Three-Scale Inputs}
\label{sec:interface}

Every task exposes a single entry point through which the model-generated
code receives host pointers for all inputs and must return host-side
outputs. The entry point is responsible for the whole pipeline---device
allocation, host-to-device transfer, kernel launch, device-to-host copy,
and deallocation---and its full wall time is measured by CUDA events on
every call. No portion of the work can be relocated into an uncounted
setup phase. This is a deliberate departure from suites that time only
the inner kernel: our target quantity is end-to-end acceleration, which
is what an applied user experiences.

To suppress cold-start noise, each evaluation runs a warmup phase
before the timed phase. Warmup repeatedly invokes the solution on the
true input until the first-call overheads---lazy CUDA context creation,
JIT compilation of any kernels dispatched through runtime-library
paths, page-locked allocation, and loading of constant memory---have
settled. Only the subsequent timed trials are aggregated into the
reported speedup. The warmup also exposes \emph{idempotence} failures:
solutions that mutate their input buffers in place, or leak device
state between calls, produce divergent trial-to-trial numbers and are
caught by the per-trial oracle rather than silently yielding an
optimistic average.

Each task ships three deterministic input scales, covering four orders
of magnitude of single-thread CPU baseline time. The \textbf{small}
scale is calibrated so that the typical task completes its CPU reference
in a handful of milliseconds---a median near ten milliseconds across the
suite, with the longest single task under six seconds---making the full
43-task sweep tractable for continuous integration on a laptop-class
GPU. The \textbf{medium} scale, used for our headline leaderboard, shifts
the central tendency of CPU time up by roughly an order of magnitude
(median under half a second, average around three seconds, worst case
about twenty seconds), so that fixed GPU overheads such as H2D transfer
and kernel-launch latency are meaningfully amortized and speedup
differences between implementations become separable. The \textbf{large}
scale is deliberately heavy-tailed: the median CPU run is a few seconds,
but the average is well over two minutes and the heaviest tasks---dense
combinatorial solvers and pentadiagonal systems---run for tens of minutes
on a single core, stressing both algorithmic scalability and memory
hierarchy utilization on the GPU side. The three-scale design lets us
observe how an LLM-generated kernel's speedup trends with problem size:
some models produce kernels that scale well (speedup grows with $N$),
while others plateau or degenerate on larger inputs.

\subsection{Correctness Oracle}
The solution's textual output is compared against the expected output
produced by the CPU reference. Each task declares one of two modes:
\textbf{exact} string comparison (integer or combinatorial outputs) or
\textbf{numerical} comparison via element-wise absolute and relative
tolerances, with task-specific $\mathrm{atol}$ and $\mathrm{rtol}$.
Tolerances are chosen tight enough to reject algorithmic bugs but loose
enough to accommodate IEEE-754 rounding under reordered floating-point
reductions on the GPU.

\subsection{Prompting Levels}

\begin{table}[ht]
\centering
\small
\caption{The three prompting regimes. Unless stated otherwise, experiments in Section~\ref{sec:results} use L3.}
\label{tab:levels}
\begin{tabular}{@{}lll@{}}
\toprule
\textbf{Level} & \textbf{Hints given} & \textbf{Purpose} \\
\midrule
L1 & Algorithm background + optimization plan & Measure ceiling with scaffolding \\
L2 & Brief optimization hints & Measure with minimal guidance \\
L3 & None (task description only) & Measure autonomous capability \\
\bottomrule
\end{tabular}
\end{table}

\subsection{Task Distribution}

\begin{table}[ht]
\centering
\small
\caption{Distribution of the 43 curated tasks across algorithmic categories.}
\label{tab:categories}
\begin{tabular}{@{}lrl@{}}
\toprule
\textbf{Category} & \textbf{\#} & \textbf{Representative tasks} \\
\midrule
Dynamic programming (OR)     & 9 & network revenue mgmt, robust value iteration, inventory replen. \\
Financial computing          & 5 & Black--Scholes, Monte-Carlo, bond and repo pricing \\
Graph analytics              & 5 & PageRank, triangle counting, connected components, BFS \\
Sparse linear algebra        & 4 & SpMV, symmetric Gauss--Seidel, CG, multigrid \\
Fluid simulation (SPH)       & 3 & cell index, force computation, position integration \\
Spatial distance / cluster.  & 4 & DBSCAN, DTW, Euclidean, Hausdorff \\
Optimization / LP            & 3 & PDLP, crew pairing, Motzkin--Straus \\
CFD / stencil                & 4 & miniWeather, LU-SSOR, ADI pentadiagonal, hotspot \\
Bioinformatics, geometry, automata, MAB & 6 & Smith--Waterman, regex, Thompson sampling, collision \\
\midrule
\textbf{Total}               & \textbf{43} & \\
\bottomrule
\end{tabular}
\end{table}

\section{Evaluation Methodology}
\label{sec:eval}

\subsection{Pipeline}
For each \((\text{model}, \text{task}, \text{sample})\) tuple:
\begin{enumerate}
  \item \textbf{Generate}: call the model with the L3 prompt; extract CUDA code from the response.
  \item \textbf{Compile}: \texttt{nvcc -O2 -arch=sm\_89} over \texttt{harness\_gpu.cu + task\_io.cu + solution.cu}.
  \item \textbf{Execute}: run with \texttt{--validate}, producing \texttt{output.txt} and \texttt{timing.json}.
  \item \textbf{Validate}: compare \texttt{output.txt} against the precomputed expected output.
  \item \textbf{Benchmark}: extract mean/std/min/max of CUDA Event timings; report end-to-end and kernel-only speedups over the pinned CPU baseline.
\end{enumerate}

\subsection{Metrics}
\label{sec:metrics}

We report four headline metrics per (model, size) cell. Let $\mathcal{T}$
denote the set of $N$ active tasks and $\mathcal{P}_m \subseteq \mathcal{T}$
the subset where model $m$'s solution \emph{compiles and produces correct
output} (``passes''). Let $s_{m,t} = T_{\mathrm{CPU},t} / T_{\mathrm{GPU},m,t}$
be the end-to-end speedup of model $m$ on task $t$ over the pinned CPU
baseline.

\begin{description}
\item[\#~pass / $N$] Cardinality of $\mathcal{P}_m$, i.e.\ the count of tasks
  the model both compiles and gets bitwise-correct output for.
  Reported alongside $N$ to avoid the misleading large-$\mathcal{P}_m$,
  small-$N$ comparison.

\item[$\mathrm{GM}\!\times$~(pass)]
  Geometric mean of speedups over \emph{passed} tasks:
  $\mathrm{GM}_m = \big(\prod_{t \in \mathcal{P}_m} s_{m,t}\big)^{1/|\mathcal{P}_m|}$.
  We use geometric -- not arithmetic -- mean because speedups are
  multiplicative ratios; arithmetic mean would be dominated by single
  outliers (e.g.\ a $2{,}500\times$ task), and would also be sensitive to
  whether one measures CPU/GPU or GPU/CPU. The geometric mean is the
  unique aggregate consistent with the multiplicative scale.

\item[$\mathrm{PAR}$-$\mathrm{GM}$~(pass)] \emph{Peak Attainment Ratio},
  geometric mean. For each task $t$ we define a per-task best
  $b_t = \max_{m'} s_{m',t}$ over the displayed model set. The model's PAR
  on task $t$ is $r_{m,t} = s_{m,t} / b_t \in [0, 1]$. We then aggregate
  $\mathrm{PAR}\text{-}\mathrm{GM}_m =
   \big(\prod_{t \in \mathcal{P}_m} r_{m,t}\big)^{1/|\mathcal{P}_m|}$.
  Intuition: a per-task $r_{m,t}$ of $1$ means the model \emph{is} the
  per-task winner; PAR-GM near $1$ means the model is consistently a
  per-task champion. Unlike raw speedups, PAR is task-difficulty-normalized:
  a task with a $2{,}500\times$ ceiling and one with a $5\times$ ceiling
  contribute symmetrically. We use this metric as the primary cross-model
  comparator and also as the headline ``replication'' metric in the
  strategy-transfer experiment (\S\ref{sec:strategy}), where $b_t$ is held
  fixed at the per-task winner of Stage~1.

\item[$\mathrm{fast}_{10}$] Fraction of $N$ tasks (\emph{not} of
  $|\mathcal{P}_m|$) for which the model achieves $s_{m,t} \geq 10$. The
  $N$-denominator is intentional: a model that passes few tasks but with
  high speedup should not appear competitive on this metric. We use $10$
  as the threshold because it is the qualitative break between
  ``GPU-as-accelerator'' ($\geq 10\times$) and ``GPU-as-rewrite-experiment''
  (lower).
\end{description}

We deliberately omit raw arithmetic means and median speedup over passed
tasks. The arithmetic mean is statistically unsound for ratios; the
median, while robust, hides single high-impact wins (e.g.\ Claude's
$696\times$ lift on \texttt{npb\_lu\_ssor\_structured} after guidance) and
is not invariant under the choice of measurement direction. Compile rate
and pass rate are not separately reported in the leaderboard because the
\#~pass / $N$ column already encodes both ($\#$~pass $\leq$ \#~compile
$\leq N$); the difference is broken out in failure-mode analyses
(\S\ref{sec:failures} and Appendix~\ref{app:badcases}).

\subsection{Setup}
All experiments run on a single NVIDIA L40S (\texttt{sm\_89}) with CUDA~12.4 and GCC~11. CPU baselines use a single thread pinned to core~0. Each model generates one sample per task at L3 with temperature~0.7.

\section{Experiments}
\label{sec:results}

\subsection{Overall Results}

Table~\ref{tab:overall} reports the seven models on L3 at the medium scale, our designated leaderboard size. Gemini~3.1~Pro leads on every metric: 83.3\% pass rate, $33\times$ median speedup, and 54.8\% of tasks crossing the $10\times$ threshold. Claude~Opus~4.6 follows at 78.6\% pass. Open-weight models (Qwen, Kimi, GLM, DeepSeek) cluster below: GLM~5.1 matches Claude in pass rate but achieves only $9.4\times$ median speedup, indicating that ``correctness'' and ``fast correctness'' are distinct skills. DeepSeek~V3.2 is the weakest model overall, with 31\% of tasks failing to compile---primarily due to missing CUDA library includes and incorrect template instantiations.

\begin{table}[ht]
\centering
\small
\caption{Overall results on \textsc{AccelEval}-43 (L3, single sample) at the \textbf{medium} scale. ``Sp.'' denotes end-to-end speedup over the pinned CPU baseline; $\mathrm{fast}_p$ is the fraction of tasks whose speedup is $\geq p\times$ (computed over all 42 evaluated tasks, so compile/correctness failures count as $0$ speedup).}
\label{tab:overall}
\begin{tabular}{@{}lrrrrrr@{}}
\toprule
\textbf{Model} & \textbf{Compile} & \textbf{Pass} & \textbf{Median Sp.} & \textbf{$\mathrm{fast}_1$} & \textbf{$\mathrm{fast}_5$} & \textbf{$\mathrm{fast}_{10}$} \\
\midrule
Gemini 3.1 Pro     & 92.9\% & \textbf{83.3\%} & \textbf{32.9$\times$} & \textbf{66.7\%} & \textbf{59.5\%} & \textbf{54.8\%} \\
Claude Opus 4.6    & \textbf{92.9\%} & 78.6\% & 16.9$\times$ & 57.1\% & 45.2\% & 40.5\% \\
GLM 5.1            & 85.7\% & 76.2\% &  9.4$\times$ & 57.1\% & 45.2\% & 38.1\% \\
Qwen 3.6 Plus      & 90.5\% & 73.8\% & 15.3$\times$ & 52.4\% & 40.5\% & 40.5\% \\
Kimi K2.5          & 78.6\% & 64.3\% & 12.0$\times$ & 45.2\% & 40.5\% & 33.3\% \\
GPT-5.4            & 81.0\% & 59.5\% & 18.7$\times$ & 42.9\% & 35.7\% & 35.7\% \\
DeepSeek V3.2      & 69.0\% & 59.5\% &  7.3$\times$ & 40.5\% & 31.0\% & 19.0\% \\
\bottomrule
\end{tabular}
\end{table}

\subsection{Scale Sensitivity}

Table~\ref{tab:scales-results} cross-tabulates median speedup at each scale. The rank order is stable, but the absolute speedup grows steeply with input size: Gemini's median speedup rises from $15\times$ (small) to $33\times$ (medium) to $93\times$ (large), reflecting the familiar GPU curve in which fixed launch/transfer overhead amortizes over larger problems. Pass rates shift less, which is consistent with the view that passing or failing is determined by \emph{algorithmic correctness} (a scale-invariant property) while speedup is dominated by \emph{resource utilization} at scale.

\begin{table}[ht]
\centering
\small
\caption{Median end-to-end speedup at three input scales (L3). Pass rate is reported in parentheses. Large-scale entries marked (---) are pending in a still-running evaluation pass (see Section~\ref{sec:limitations}).}
\label{tab:scales-results}
\begin{tabular}{@{}lrrr@{}}
\toprule
\textbf{Model} & \textbf{Small} & \textbf{Medium} & \textbf{Large} \\
\midrule
Gemini 3.1 Pro    & 15.1$\times$ (83.3\%) & 32.9$\times$ (83.3\%) & \textbf{93.1$\times$ (85.7\%)} \\
Claude Opus 4.6   &  7.6$\times$ (81.0\%) & 16.9$\times$ (78.6\%) & 28.3$\times$ (76.2\%) \\
GPT-5.4           &  4.1$\times$ (59.5\%) & 18.7$\times$ (59.5\%) & --- \\
GLM 5.1           &  7.9$\times$ (76.2\%) &  9.4$\times$ (76.2\%) & --- \\
Qwen 3.6 Plus     &  2.4$\times$ (76.2\%) & 15.3$\times$ (73.8\%) & --- \\
Kimi K2.5         &  5.5$\times$ (71.4\%) & 12.0$\times$ (64.3\%) & --- \\
DeepSeek V3.2     &  8.2$\times$ (64.3\%) &  7.3$\times$ (59.5\%) & --- \\
\bottomrule
\end{tabular}
\end{table}

\subsection{Per-Category Analysis}

Category-level pass rates are heterogeneous:
\begin{itemize}
  \item \textbf{Embarrassingly parallel (Monte~Carlo, Black--Scholes, element-wise stencils)}: all seven models pass consistently with $>50\times$ speedup on medium. Here the L3 gap between frontier and open-weight is narrow.
  \item \textbf{Dynamic programming (Bellman-Ford, held-Karp TSP, HJB value iteration, inventory-replenishment)}: Gemini and Claude lead; open-weight models frequently produce correct but serial ``one-block'' kernels that underutilize the GPU.
  \item \textbf{Sparse linear algebra (HPCG SpMV/SymGS/MG, NPB CG, SpMV-CSR)}: the sharpest differentiator. Sparse patterns defeat the four weakest models on multiple tasks; only Gemini attains consistent correctness with meaningful speedup on the full sparse suite.
  \item \textbf{Graph analytics (GAP PageRank, triangle count, connected components, BFS, push-relabel max-flow)}: Claude edges out Gemini on push-relabel; Gemini leads on the rest.
  \item \textbf{Financial computing (bonds pricing, repo pricing, batched portfolio optimization)}: all leaders use cuBLAS/cuSOLVER; models that attempt hand-written LU factorization typically produce numerically unstable results at medium scale.
\end{itemize}

\subsection{Failure Modes}
\label{sec:failures}

We classify failures into four buckets, counted across the $7 \times 42 = 294$ medium evaluations:
\begin{itemize}
  \item \textbf{Compile fail} ($\approx 14\%$): missing or mis-spelled CUDA library headers (cusolverDn*, cub::*), incorrect kernel signature against the declared interface, invalid template instantiation. DeepSeek~V3.2 accounts for $30\%$ of all compile failures.
  \item \textbf{Correctness fail} ($\approx 12\%$): compiled and ran, but output mismatched the oracle. Root causes concentrate in race conditions on reductions, off-by-one errors at tile boundaries, and incorrect handling of the in-place update pattern in Gauss-Seidel sweeps.
  \item \textbf{Correct but slow} ($\approx 6\%$): $<1\times$ speedup. Almost always excessive synchronization or the aforementioned ``serial one-block'' anti-pattern.
  \item \textbf{Universal failures}: on two tasks (\texttt{hpcg\_mg\_vcycle} multigrid V-cycle, \texttt{npb\_sp\_adi\_pentadiagonal} ADI pentadiagonal solve), fewer than three of the seven models produce a correct solution at medium scale. Both involve cross-iteration data dependencies that resist naive kernel extraction.
\end{itemize}

\section{Decomposition Analysis: What Makes the Kernel Fast?}
\label{sec:decomp}

Leaderboard metrics say \emph{which} model wins; they do not say \emph{why}. A 300$\times$ speedup could be driven by shared-memory tiling, algorithmic substitution, or trivial precomputation. To attribute performance to concrete techniques we build a two-stage pipeline on top of \textsc{AccelEval}: an \textbf{optimization pattern database} (Section~\ref{sec:pat-db}) and a \textbf{pairwise diff decomposition} over multi-turn agent runs (Section~\ref{sec:pairwise}).

\subsection{Optimization Pattern Database}
\label{sec:pat-db}

\begin{figure}[ht]
\centering
% \includegraphics[width=\linewidth]{figs/pattern_db.pdf}
\caption{Three-stage pipeline for the pattern database. Auto-detect uses static CUDA analysis (\texttt{\_\_shared\_\_}, \texttt{\_\_constant\_\_}, \texttt{template<>}, \texttt{\#pragma unroll}, \texttt{atomicAdd}, \texttt{\_\_shfl\_sync}, \ldots); an LLM analyzer reads the CPU reference and generated CUDA to summarize detected patterns and propose new candidates; the knowledge base accumulates evidence across runs/models, with auto-promotion once a candidate appears in sufficiently many passing samples.}
\label{fig:patdb}
\end{figure}

Each sample that compiles and passes validation is parsed by two complementary analyzers:
(i) a \emph{static detector} scans for canonical CUDA idioms (shared-memory usage, warp intrinsics, template dispatch, atomic ops, restricted pointers, \texttt{\#pragma unroll}, etc.); and
(ii) an \emph{LLM analyzer} (Gemini~3.1~Pro, budget-bounded) reads both the CPU reference and the generated CUDA and produces a natural-language pattern list, including novel candidates not in the registry.
Detected patterns are written into a pattern database with evidence counts and speedup annotations. Candidates are auto-promoted to first-class patterns once they accumulate evidence across multiple models and tasks.

\begin{table}[ht]
\centering
\small
\caption{Example patterns in the database (selected).}
\label{tab:patterns}
\begin{tabular}{@{}llrl@{}}
\toprule
\textbf{ID} & \textbf{Pattern} & \textbf{Evidence} & \textbf{Observed lift} \\
\midrule
PAT-001 & Shared-memory tiling                & 42 & $12.3\times$ avg lift \\
PAT-005 & Bitmask state representation        & 8  & $3{,}186\times$ (NFA) \\
PAT-008 & Template kernel dispatch            & 15 & $8.7\times$ avg lift \\
PAT-012 & Precomputation ($\varepsilon$-closure fusion) & 6  & eliminates BFS phase \\
\bottomrule
\end{tabular}
\end{table}

Across the 40 tasks that produced a passing solution, our agent discovered \textbf{33 distinct optimization strategies}, with the most frequent being \emph{persistent device-memory allocation} (35 tasks), \emph{precomputation via init kernel} (32 tasks), and \emph{memory coalescing} (31 tasks). The long tail contains task-specific techniques (\emph{e.g.}, $\varepsilon$-closure fusion for regex, grid-stride loops for reductions) that appear in only one or two tasks. A strategy-intensity matrix (tasks~$\times$~patterns) reveals which categories are ``saturated'' by a few common patterns and which require bespoke strategies.

\subsection{Pairwise Diff Attribution}
\label{sec:pairwise}

\textsc{AccelEval} supports a \emph{multi-turn agent} mode in which the model is given feedback (compilation errors, correctness mismatches, slow kernel profile) and produces successive versions $V_0, V_1, \ldots, V_k$ of the same solution. Consecutive versions form natural A/B experiments: the same task, same oracle, and two code snapshots differing by a small edit. We attribute speedup change to the concrete patterns that changed between them.

Formally, for each consecutive pair $(V_a, V_b)$ we record a \texttt{DiffRecord}:
\begin{itemize}
  \item \texttt{pattern\_changes[]}: patterns added/removed/modified between $V_a$ and $V_b$;
  \item \texttt{causal\_chains[]}: ordered sequences inferred across multiple diffs (\emph{e.g.}, shared-memory tiling enables subsequent template specialization);
  \item \texttt{new\_candidates[]}: previously unseen techniques flagged by the LLM analyzer;
  \item \texttt{speedup\_ratio} $= T(V_a) / T(V_b)$.
\end{itemize}

Aggregating \texttt{DiffRecord}s across all multi-turn traces yields a \emph{pattern$\to$speedup attribution table}. In one representative trace ($V_0 = 5.2\times$, $V_4 = 48.7\times$), the decomposition attributes $4.1\times$ of the $9.4\times$ end-to-end improvement to shared-memory tiling, $2.3\times$ to template dispatch, and the remainder to bounds-check elimination.

\paragraph{Why pairwise (not absolute).}
Absolute correlations between ``pattern present'' and ``speedup'' are confounded: fast kernels often bundle many patterns. Pairwise diffs isolate the marginal contribution of each edit, approximating the counterfactual ``what if only this pattern changed?'' without requiring manual ablations.

\paragraph{Limitations.}
Attribution is only as reliable as the granularity of the agent's edits. Large, multi-pattern diffs provide weaker causal claims than small ones. Our current pipeline reports both per-pattern lifts and bundle-level lifts for transparency.

\section{Strategy Recommendation: Three Baselines}
\label{sec:strategy}

Using the formulation of Section~\ref{sec:formulation}, we benchmark
three baselines of increasing complexity. All three share the same
feature extractor $f$ and feasible set $\mathcal{A}_0$; they differ
only in how they estimate the effect function $\hat{g}$ in
Equation~\eqref{eq:objective}.

\subsection{Baseline 1: Nearest-Neighbor Transfer}
\label{sec:bl1}
Retrieve the most similar problem in the history and reuse its
best-observed strategy:
\begin{equation}
j^* = \arg\min_{j}\;\bigl\lVert f(Q^*) - f(Q_j)\bigr\rVert_2,
\qquad
\mathbf{s}^* = \arg\max_{\substack{\mathbf{s}:\;\exists\,l \text{ s.t. } j(l)=j^*,\; \mathbf{s}_l^B = \mathbf{s}}} r_l.
\label{eq:nn}
\end{equation}
We also report the $k$-NN variant with distance-weighted voting. This
baseline learns no model; it relies purely on the locality assumption
in feature space.

\subsection{Baseline 2: Greedy Coordinate Search}
\label{sec:bl2}
Starting from $\mathbf{s} = \mathbf{s}^A$, sweep one dimension at a
time. In round $i$, hold all other coordinates fixed and choose the
intensity that maximizes observed speedup:
\begin{equation}
s_i^* = \arg\max_{s_i \in \{0,\ldots,M_i\}}\,
  y\bigl(\mathbf{s}[s_i \!\to\! s_i],\; Q^*\bigr),
  \qquad i = 1, \ldots, m.
\label{eq:greedy}
\end{equation}
Total evaluations scale as $\sum_{i=1}^{m}(M_i + 1)$, compared to the
exhaustive $\prod_i(M_i+1)$. Greedy search ignores cross-dimension
interactions but costs only linear evaluations.

\subsection{Baseline 3: Linear Regression}
\label{sec:bl3}
Fit the effect function as a linear model in features and strategy:
\begin{equation}
\hat{g}(\mathbf{q}, \mathbf{s}^A, \Delta\mathbf{s})
  = \mathbf{w}_q^{\top}\mathbf{q}
  + \mathbf{w}_s^{\top}\mathbf{s}^A
  + \mathbf{w}_\Delta^{\top}\Delta\mathbf{s} + b,
\label{eq:linreg}
\end{equation}
learned by least squares on $\{(\mathbf{x}_l, r_l)\}_{l=1}^{k}$.
Substituting into Equation~\eqref{eq:objective} gives a separable
optimum per coordinate (or an integer-linear program when
$\mathcal{A}_0$ is binding). Linear regression cannot capture
interactions but provides an interpretable, principled lower bound.

\subsection{Experimental Protocol}

\paragraph{Data split.} We hold out $20\%$ of $\mathcal{P}$ as
unseen queries (``test tasks'') and train all three baselines on the
remaining $80\%$, using the multi-turn agent traces from
Section~\ref{sec:decomp} as the history $\{E_l\}$.

\paragraph{Metrics.}
\begin{itemize}
  \item \textbf{Hit@$K$} --- fraction of test queries whose predicted
  $\mathbf{s}^*$ (top-$K$) matches the best-observed strategy in the
  held-out trace.
  \item \textbf{Speedup gap} --- $\Delta = y(\mathbf{s}^{\mathrm{oracle}})
  - y(\mathbf{s}^*)$, averaged over test queries.
  \item \textbf{Evaluation budget} --- number of compile-and-run trials
  needed to reach within $10\%$ of the oracle speedup.
\end{itemize}

\subsection{Results}
\label{sec:strategy-results}

\begin{table}[ht]
\centering
\small
\caption{Strategy-recommendation baselines on \textsc{AccelEval}-42.
\emph{(Results to be filled.)} $\mathbf{s}^A = \mathbf{0}$ (raw CPU);
history is drawn from the multi-turn traces of
Section~\ref{sec:decomp}.}
\label{tab:strategy}
\begin{tabular}{@{}lrrr@{}}
\toprule
\textbf{Method} & \textbf{Hit@1} & \textbf{Avg.\ speedup gap} & \textbf{\# trials to 0.9$\times$oracle} \\
\midrule
Baseline 1: 1-NN Transfer       & --- & --- & --- \\
Baseline 1$^\prime$: 5-NN Voting & --- & --- & --- \\
Baseline 2: Greedy Coordinate    & --- & --- & --- \\
Baseline 3: Linear Regression    & --- & --- & --- \\
\midrule
LLM (Gemini 3.1 Pro, L3)          & --- & --- & --- \\
LLM + history (retrieval-augmented) & --- & --- & --- \\
\bottomrule
\end{tabular}
\end{table}

\begin{figure}[ht]
\centering
% \includegraphics[width=0.7\linewidth]{figs/strategy_budget.pdf}
\caption{\emph{(To be filled.)} Speedup attained versus
compile-and-run budget for the three baselines and the LLM.
Expected trend: Baseline~1 saturates early at its retrieval ceiling;
Baseline~2 improves monotonically but plateaus before the oracle;
Baseline~3 is competitive at very low budgets but asymptotes below
the interaction-aware methods.}
\label{fig:strategy-budget}
\end{figure}

\paragraph{What we expect.}
We anticipate: (i)~NN-transfer will dominate on tasks with close
neighbors in $\mathcal{F}$ (graph analytics, dense linear algebra),
but collapse on rare categories (regex, combinatorial optimization);
(ii)~greedy coordinate search will outperform NN-transfer given
a modest budget, but miss interactions such as tiling$\times$unroll;
(iii)~linear regression will be the weakest on raw speedup but most
sample-efficient; (iv)~strong LLMs will roughly match or exceed the
classical baselines on tasks close to pretraining data, and fall
sharply on bespoke categories.

\section{Strategy Transfer: Discovery vs.\ Implementation}
\label{sec:transfer}

The L3 results in Section~\ref{sec:results} conflate two skills:
\emph{knowing what to optimize} (strategy discovery) and
\emph{knowing how to code it} (implementation). A model could
underperform because it cannot identify useful patterns, or because
it identifies them but writes buggy/slow CUDA. We design a two-stage
ablation to tell these apart.

\paragraph{Stage~1 (baseline, from Section~\ref{sec:results}).}
Each model generates a solution under L3. Let $M^{*}(Q)$ denote the
best model on task $Q$ (highest correct speedup).

\paragraph{Stage~2 (strategy transfer).}
For each non-winning model $M$ on task $Q$, we take $M^{*}(Q)$'s
solution, extract its optimization strategy as a natural-language
recipe (ordered list of patterns + rationale), and re-prompt $M$ with
this recipe appended to the L3 prompt. $M$ must then produce its own
implementation guided by the recipe, without seeing $M^{*}$'s code.

\paragraph{Diagnosis.} Let $y^{\mathrm{S1}}_{M}$ and
$y^{\mathrm{S2}}_{M}$ be the speedups in Stage~1 and Stage~2. Define
the \emph{replication ratio}
\begin{equation}
\mathrm{RR}(M, Q) = \frac{y^{\mathrm{S2}}_{M}(Q)}{y(M^*)(Q)}
  \in [0, 1].
\label{eq:rr}
\end{equation}
We read the outcome per task:
\begin{itemize}
  \item $\mathrm{RR} \geq \tau$ (large, close to~1): the Stage~1 gap
        was \textbf{strategy discovery} --- given the recipe, $M$
        can execute.
  \item $\mathrm{RR} < \tau$ (small): the gap is \textbf{implementation
        skill} --- even told what to do, $M$ cannot write the code.
\end{itemize}
We use $\tau = 0.7$ (within $30\%$ of the leader) as the default.

\paragraph{Why this matters.}
The two diagnoses call for different remedies. A discovery gap points
to retrieval augmentation (Section~\ref{sec:strategy}) or agentic
reasoning. An implementation gap is harder: it points to limitations
in the base model's CUDA fluency, which supervised fine-tuning or
longer context with canonical examples may fix.

\paragraph{Experimental protocol.}
\begin{enumerate}
  \item Select the $T$ tasks where Gemini~3.1~Pro (the Stage~1
        leader) achieves $\geq 10\times$ speedup and no other model is
        within $2\times$. ($T \approx 15$ in our current data.)
  \item For each competitor $M \in \{$Claude, Qwen, Kimi, GLM,
        DeepSeek$\}$, run the transfer protocol.
  \item Report per-model, per-task RR and the fraction of tasks
        classified as \emph{discovery} vs.\ \emph{implementation}.
\end{enumerate}

\begin{table}[ht]
\centering
\small
\caption{Strategy-transfer diagnosis per competitor model.
\emph{(Results to be filled.)} For each model we report the mean
replication ratio $\mathrm{RR}$, the fraction of tasks with
$\mathrm{RR} \geq 0.7$ (``discovery gap''), and the fraction with
$\mathrm{RR} < 0.7$ (``implementation gap'').}
\label{tab:transfer}
\begin{tabular}{@{}lrrrr@{}}
\toprule
\textbf{Competitor} & \textbf{\# tasks} & \textbf{Mean RR} & \textbf{Discovery\%} & \textbf{Impl.\%} \\
\midrule
Claude Opus 4.6 & --- & --- & --- & --- \\
Qwen 3.6 Plus    & --- & --- & --- & --- \\
Kimi K2.5        & --- & --- & --- & --- \\
GLM 5.1          & --- & --- & --- & --- \\
DeepSeek V3.2    & --- & --- & --- & --- \\
\bottomrule
\end{tabular}
\end{table}

\begin{figure}[ht]
\centering
% \includegraphics[width=0.55\linewidth]{figs/transfer_diagnosis.pdf}
\caption{\emph{(To be filled.)} Scatter of $y^{\mathrm{S1}}$ (x-axis)
vs.\ $y^{\mathrm{S2}}$ (y-axis) per task. Points on the $y = y^{\mathrm{S2}}_{\mathrm{leader}}$ line indicate a fully-closed discovery gap;
points clustered near the $y = y^{\mathrm{S1}}$ line indicate a
stubborn implementation gap.}
\label{fig:transfer}
\end{figure}

\section{Discussion}
\label{sec:discussion}

\paragraph{Where LLMs excel.} Models consistently produce strong CUDA code for massively parallel, regular workloads such as Monte~Carlo simulation and dense element-wise kernels, often exceeding $1000\times$ speedup.

\paragraph{Where they fail.} Tasks requiring irregular memory access patterns (sparse solvers, graph irregularity) or cross-iteration dependencies (multigrid V-cycles, Gauss-Seidel sweeps) remain fragile. Three tasks (\texttt{nbnxm\_forces}, \texttt{npb\_mg\_vcycle}, \texttt{rodinia\_nw\_alignment}) were failed by \emph{all} six models.

\paragraph{Infrastructure matters.} Routing Chinese models through OpenRouter introduces significant instability: $\sim$18\% of GLM requests returned empty responses in our runs. We added empty-response retries and throughput-based routing to partially mitigate this.

\section{Limitations}
\label{sec:limitations}

\begin{itemize}
  \item \textbf{Single sample per task.} We report one generation per $(\text{model},\text{task})$ pair; $k$-sample studies would tighten pass-rate confidence intervals.
  \item \textbf{Single GPU architecture.} All measurements on \texttt{sm\_89}; generalization to H100 / A100 / consumer cards is unstudied.
  \item \textbf{L3 only.} We have not yet scaled the L1/L2 evaluation across all models.
\end{itemize}

\section{Conclusion}
\label{sec:conclusion}

\textsc{AccelEval} provides a rigorous, reproducible platform for evaluating LLM-generated GPU code on realistic workloads. Our first large-scale study shows a clear capability gap between frontier closed-source models and open-weight alternatives, concentrated in correctness on irregular HPC kernels. We release the benchmark, the harness, and the evaluation scripts to support community benchmarking.

\section*{Broader Impact}

\textsc{AccelEval} aims to make LLM-assisted GPU programming more rigorous and reproducible. Positive impact: lowering the barrier to entry for GPU acceleration in scientific computing; enabling fair cross-model evaluation; surfacing failure modes that model developers can target. Potential concerns: benchmark saturation incentivizing overfitting; uneven access to GPU infrastructure may advantage models that produce simpler kernels. We release all task definitions, harnesses, and evaluation logs to enable independent verification.



\begin{ack}
We thank \emph{(TBD)}. No external funding to declare. No competing interests.
\end{ack}

\section*{References}

\medskip
{\small

\bibitem[Ouyang et al.(2025)]{kernelbench}
A.\ Ouyang, S.\ Guo, S.\ Arora, A.\ Zhang, M.\ Hu, C.\ R\'e, and A.\ Mirhoseini.
(2025) KernelBench: Can LLMs write efficient GPU kernels?
{\it arXiv preprint arXiv:2502.10517}.

\bibitem[Chen et al.(2021)]{humaneval}
M.\ Chen et al.\ (2021) Evaluating large language models trained on code.
{\it arXiv preprint arXiv:2107.03374}.

\bibitem[Austin et al.(2021)]{mbpp}
J.\ Austin et al.\ (2021) Program synthesis with large language models.
{\it arXiv preprint arXiv:2108.07732}.

\bibitem[Jimenez et al.(2024)]{swebench}
C.\ E.\ Jimenez et al.\ (2024) SWE-bench: Can language models resolve
real-world GitHub issues?
{\it ICLR}.

\bibitem[Wen et al.(2022)]{babeltower}
Y.\ Wen et al.\ (2022) BabelTower: Learning to auto-parallelized program
translation.
{\it ICML}.

\bibitem[Tehranijamsaz et al.(2024)]{coderosetta}
A.\ Tehranijamsaz et al.\ (2024) CodeRosetta: Pushing the boundaries of
unsupervised code translation for parallel programming.
{\it NeurIPS}.

\bibitem[Nugteren \& Corporaal(2014)]{bones}
C.\ Nugteren and H.\ Corporaal.\ (2014) Bones: An automatic skeleton-based
C-to-CUDA compiler for GPUs.
{\it ACM TACO}, 11(4).

\bibitem[Amini et al.(2012)]{par4all}
M.\ Amini et al.\ (2012) Par4All: From convex array regions to heterogeneous
computing.
{\it IMPACT}.

\bibitem[Verdoolaege et al.(2013)]{ppcg}
S.\ Verdoolaege, J.\ Carlos Juega, A.\ Cohen, J.\ I.\ G\'omez, C.\ Tenllado,
and F.\ Catthoor.\ (2013) Polyhedral parallel code generation for CUDA.
{\it ACM TACO}, 9(4):54.

\bibitem[Jain et al.(2024)]{livecodebench}
N.\ Jain et al.\ (2024) LiveCodeBench: Holistic and contamination-free evaluation
of large language models for code.
{\it arXiv preprint arXiv:2403.07974}.

\bibitem[Zhuo et al.(2024)]{bigcodebench}
T.\ Y.\ Zhuo et al.\ (2024) BigCodeBench: Benchmarking code generation with
diverse function calls and complex instructions.
{\it arXiv preprint arXiv:2406.15877}.

\bibitem[NVIDIA(2024)]{computeeval}
NVIDIA.\ (2024) ComputeEval: Evaluating LLMs on accelerated computing code
generation. \url{https://github.com/NVIDIA/compute-eval}.

\bibitem[Nichols et al.(2024)]{pareval}
D.\ Nichols, J.\ H.\ Davis, Z.\ Xie, A.\ Rajaram, and A.\ Bhatele.\ (2024)
Can large language models write parallel code?
{\it HPDC}.

\bibitem[Li et al.(2025)]{tritonbench}
Z.\ Li et al.\ (2025) TritonBench: Benchmarking large language model
capabilities for generating Triton operators.
{\it arXiv preprint arXiv:2502.14752}.

\bibitem[Wei et al.(2025)]{multikernelbench}
Z.\ Wei et al.\ (2025) MultiKernelBench: A multi-platform benchmark for
kernel generation.
{\it arXiv preprint}.

\bibitem[Beamer et al.(2015)]{gapbs}
S.\ Beamer, K.\ Asanovi\'c, and D.\ Patterson.\ (2015) The GAP Benchmark Suite.
{\it arXiv preprint arXiv:1508.03619}.

\bibitem[Dongarra et al.(2016)]{hpcg}
J.\ Dongarra, M.\ A.\ Heroux, and P.\ Luszczek.\ (2016) High-performance conjugate-gradient
benchmark: A new metric for ranking high-performance computing systems.
{\it IJHPCA}, 30(1):3--10.

\bibitem[Bailey et al.(1991)]{npb}
D.\ H.\ Bailey et al.\ (1991) The NAS Parallel Benchmarks.
{\it IJHPCA}, 5(3):63--73.

\bibitem[Che et al.(2009)]{rodinia}
S.\ Che et al.\ (2009) Rodinia: A benchmark suite for heterogeneous computing.
{\it IISWC}.

\bibitem[Norman et al.(2020)]{miniweather}
M.\ Norman et al.\ (2020) miniWeather: A mini-app for teaching parallel programming.
\url{https://github.com/mrnorman/miniWeather}.

\bibitem[van Werkhoven(2019)]{kerneltuner}
B.\ van Werkhoven.\ (2019) Kernel Tuner: A search-optimizing GPU code auto-tuner.
{\it FGCS}, 90:347--358.

\bibitem[Ansel et al.(2014)]{opentuner}
J.\ Ansel et al.\ (2014) OpenTuner: An extensible framework for program autotuning.
{\it PACT}.

\bibitem[Whaley \& Dongarra(1998)]{atlas}
R.\ C.\ Whaley and J.\ J.\ Dongarra.\ (1998) Automatically tuned linear algebra software.
{\it SC}.

\bibitem[Adams et al.(2019)]{halide-autoscheduler}
A.\ Adams et al.\ (2019) Learning to optimize halide with tree search and
random programs.
{\it ACM TOG}, 38(4).

\bibitem[Chen et al.(2018)]{tvm}
T.\ Chen et al.\ (2018) TVM: An automated end-to-end optimizing compiler for
deep learning.
{\it OSDI}.
}

%%%%%%%%%%%%%%%%%%%%%%%%%%%%%%%%%%%%%%%%%%%%%%%%%%%%%%%%%%%%

\appendix

\section{Complete Task List}
\label{app:tasklist}

\emph{(42 tasks with category, difficulty, and brief description.
Auto-generated from \texttt{tasks/*/task.json}.)}

\section{Example Prompt (L3)}
\label{app:prompt}

\begin{lstlisting}[style=cuda,caption={Canonical L3 prompt structure.}]
# Task: <name>

<task description>

## Interface
<C function signatures>

## Input sizes
<auto-generated size table>

## Output contract
<expected output format>

## CPU Reference
<cpu_reference.c content>

Please produce a single self-contained CUDA file that implements
<function names>. Your code will be compiled and linked against
the provided harness.
\end{lstlisting}

%%%%%%%%%%%%%%%%%%%%%%%%%%%%%%%%%%%%%%%%%%%%%%%%%%%%%%%%%%%%

\section{Bad-Case Analysis: Three Representative Failure Patterns}
\label{app:badcases}

This appendix dissects three tasks with high failure rates across
the eight evaluated models. Each illustrates a distinct failure mode
that recurs throughout the benchmark: numerical fragility, an
intrinsic GPU ceiling, and a toolchain limitation. Together they
explain $\approx 35\%$ of all observed correctness or compile failures
on \texttt{medium} inputs.

\subsection{Numerical Fragility: \texttt{bonds\_pricing}}
\label{app:badcase-bonds}

\textbf{Task.} Faithful port of FinanceBench bond-pricing kernels --
embarrassingly parallel over $N$ bonds, but each bond requires
yield-to-maturity Newton iteration plus discounted-cashflow
summation, where the reference uses \texttt{double} throughout.

\textbf{Outcome.} Of eight models, only one (\textsc{Gemini}) passes
\texttt{medium}, achieving $120.6\times$. The remaining seven split
into two compile failures and five correctness failures. Despite the
algorithm being trivially parallel, every correctness failure
manifests as a different first-bond price, exposing distinct bugs:

\begin{table}[h]
\centering
\small
\begin{tabular}{lrrrl}
\toprule
Model & First fail & Got & Expected & Root cause \\
\midrule
\textsc{GPT-5.4}  & idx 0   & $1{\times}10^{30}$ & $93.45$ & sentinel write-through \\
\textsc{V4 Pro}   & idx 3   & $8.4{\times}10^{26}$ & $1.9{\times}10^{8}$ & \texttt{MAX\_LEGS=32} cap \\
\textsc{Qwen 3.6} & idx 0   & $117.4$ & $93.45$ & all-\texttt{float} Newton diverges \\
\textsc{Kimi K2.5}& idx 0   & $103.4$ & $93.45$ & cashflow off-by-one \\
\textsc{V3.2}     & idx 116 & $93.53$ & $93.51$ & \texttt{(float)} cast at output \\
\bottomrule
\end{tabular}
\caption{Per-model correctness failure on \texttt{bonds\_pricing}.}
\end{table}

\textit{Pattern A -- defensive sentinel.} \textsc{GPT-5.4} introduces
\texttt{\#define INVALID\_RESULT 1.0e30f} and writes it whenever a
self-imposed input check fails:
\begin{lstlisting}[style=cuda,basicstyle=\ttfamily\footnotesize]
if (im < 1 || im > 12 || mm < 1 || mm > 12 ||
    id < 1 || id > 31 || md < 1 || md > 31 ||
    coupon_freq_decoded <= 0.0f) {
    prices[i*4 + 0] = INVALID_RESULT; ... return;  // fires for every bond
}
\end{lstlisting}
The \texttt{coupon\_freq} decoder the model wrote is the bug; the
guard fires for every input. \textsc{V4 Pro} hard-codes
\texttt{MAX\_LEGS = 32} and writes a 1e30 sentinel when a bond
exceeds it; medium-scale 30-year semi-annual bonds have 60 legs,
tripping the cap from the third bond onward. The sentinel then
propagates through downstream arithmetic, yielding $8{\times}10^{26}$.

\textit{Pattern B -- precision-induced divergence.} \textsc{Qwen 3.6}
performs the entire Newton iteration in single precision:
\begin{lstlisting}[style=cuda,basicstyle=\ttfamily\footnotesize]
float dy = -f_val / dnpv;  y += dy;          // catastrophic cancellation
npv_final += coupons[k] * powf(1.0f + y/freq, -freq * t);
\end{lstlisting}
The Jacobian becomes ill-conditioned and the iterate latches at the
clipped upper bound $y=0.5$. The resulting price is $26\%$ too high.

\textit{Pattern C -- structural off-by-one.} \textsc{Kimi K2.5} treats
\texttt{numCashFlows} as both ``coupon count'' and ``coupon plus
redemption count,'' overwriting the last coupon leg with the
redemption payment. The error is one missing coupon's present value
($\approx 10\%$).

\textit{Pattern D -- output-side truncation.} \textsc{V3.2} computes
in \texttt{double} but stores via \texttt{prices[i] = (float)dirty}.
By the 116th bond, accumulated rounding crosses
$\mathrm{rtol}=10^{-4}$. This is the \emph{only} failure within
$0.1\%$ of the reference and would pass under either a wider
tolerance or \texttt{double} output buffers.

\textit{Compile failures.} \textsc{Claude Opus 4.6} emitted a literal
\texttt{```c} markdown fence into the \texttt{.cu} file (output
formatting bug). \textsc{GLM 5.1} wrote a static initializer with
more entries than the declared array dimension.

\textbf{Takeaway.} The five correctness failures are not a single
``precision issue'' but four orthogonal bugs: defensive sentinels
(\textsc{GPT-5.4}, \textsc{V4 Pro}), wrong-precision algorithm
(\textsc{Qwen}), off-by-one structural mistake (\textsc{Kimi}), and
final-stage truncation (\textsc{V3.2}). A benchmark that reports only
``correct/incorrect'' obscures this diversity.

\subsection{Intrinsic GPU Ceiling: \texttt{hpcg\_mg\_vcycle}}
\label{app:badcase-hpcg}

\textbf{Task.} HPCG geometric multigrid V-cycle: pre-smooth, restrict
residual, recurse on coarser grid, prolongate correction,
post-smooth. The CPU reference at \texttt{medium} runs in
$180\,\mathrm{ms}$.

\textbf{Outcome.} Only \textsc{Gemini} produces a correct solution,
and it runs at $0.70\times$ -- \emph{slower than the CPU baseline}.
The other seven models split: two compile failures
(\textsc{Kimi}: ill-formed \texttt{cudaMalloc} macro;
\textsc{V3.2}: chained pointer assignment across types) and five
correctness failures.

\textbf{Why even the winner is sub-CPU.} The \texttt{nsys} profile
of Gemini's solution shows
\texttt{memcpy\_overhead\_ms} $=682\,\mathrm{ms}$ against
\texttt{cpu\_baseline\_ms} $=180\,\mathrm{ms}$. The host-to-device
transfer alone is nearly $4\times$ the entire CPU runtime; multigrid
on a $128^3$-class grid does not amortize the transfer when the
arithmetic intensity is bounded by sparse 27-point stencil
operations. We retain this task in the suite as a deliberate
\emph{negative-result anchor}: any benchmark that does not contain at
least one task where the GPU loses must be suspect.

\textbf{Why the non-winners fail.} The five correctness failures all
involve mishandled grid coarsening: most models compute the
fine-to-coarse mapping (\texttt{f2c}) inline rather than precomputing
it, leading to off-by-one errors at level boundaries that compound
multiplicatively across V-cycles. Multigrid is famously the
``research-grade'' kernel of HPC; its inclusion is intentional.

\subsection{Toolchain Limitation: \texttt{batched\_lhpca\_portfolio}}
\label{app:badcase-lhpca}

\textbf{Task.} Batched portfolio optimization requiring repeated
Cholesky factorization of $N{\times}N$ covariance matrices and
matrix-vector products. The CPU reference uses
\texttt{malloc}-allocated dense linear algebra; CPU \texttt{medium}
takes $19.8\,\mathrm{s}$.

\textbf{Outcome.} Only two models pass: \textsc{Gemini} ($103.7\times$,
self-rolled batched Cholesky) and \textsc{Qwen} ($1.9\times$,
brute-force without library calls). The remaining six all fail at
build time, four with linker errors of the form
\texttt{undefined reference to `cublasCreate\_v2'} and
\texttt{`cusolverDnSpotrf\_bufferSize'}.

\textbf{Root cause is the toolchain, not the model.} The benchmark's
default build command (\texttt{compile.py}) is
\begin{lstlisting}[basicstyle=\ttfamily\footnotesize]
nvcc -O2 -arch=sm_89 -I include harness.cu task_io.cu solution.cu -o solution
\end{lstlisting}
with no \texttt{-lcublas -lcusolver}. Models that reach for
\textsc{cuSOLVER} for batched Cholesky -- the canonical correct choice
for this workload -- are punished by the link step. Self-implemented
Cholesky (\textsc{Gemini}) wins; brute-force (\textsc{Qwen}) compiles
but yields only $1.9\times$ because it cannot exploit
warp-cooperative triangular solves at the scale a tuned library
would.

\textbf{Takeaway.} This task exposes a benchmark-honesty issue: a
compilation-only environment under-rewards the most idiomatic CUDA
solution. We document this rather than ``fix'' it because the
$103.7\times$ ceiling Gemini achieves \emph{without} libraries
demonstrates the algorithmic headroom available, and the
two-pass design (with and without libraries) would require
re-instrumenting every task. An informed reader of the leaderboard
should view per-task best as the meaningful comparison and treat
\texttt{batched\_lhpca\_portfolio} as a measure of
``CUDA-from-scratch'' rather than ``CUDA-with-ecosystem.''

\subsection{Cross-Cutting Observation}

The three patterns are mutually disjoint. Numerical fragility
(\S\ref{app:badcase-bonds}) is a model property; the GPU ceiling
(\S\ref{app:badcase-hpcg}) is a workload property; the toolchain
limitation (\S\ref{app:badcase-lhpca}) is a benchmark-design property.
Strategy-transfer experiments
(\S\ref{sec:strategy} in the main text) are most effective for the
first category, ineffective for the second, and orthogonal to the
third.

%%%%%%%%%%%%%%%%%%%%%%%%%%%%%%%%%%%%%%%%%%%%%%%%%%%%%%%%%%%%%

\newpage
\input{checklist.tex}


\end{document}